\begin{abstract}
Inpainting models fill a hidden region of an image using its visible surroundings, and their fills are usually compared pixel-by-pixel with the original content.
A low error can mean that the fill recovered this person's own detail, or only that it looks like an average person.
For a model output, the right score is unknown, so the metric itself cannot be checked.
We therefore audit the metrics on references whose relationship to the hidden region is known in advance: a second image of the same person, the opposite side of the same body, unrelated people, and averages of many people.
We add three constructed controls: the person's own second image, pulled toward the dataset average; the same amount of detail taken from a random person; and white noise.
We apply this identity ladder to nine image sets of brain magnetic resonance imaging (MRI), retinal photographs, faces, and irises.
On a scale where 0 is the target and 1 is a random person, with the hidden region 1/8 as wide as the imaged structure, mean squared error (MSE) places the person's second image at 0.41 and the dataset average at 0.50.
The second image pulled toward the average scores better (0.29) than the unchanged one.
Fr\'echet inception distance (FID) and Fr\'echet wavelet distance (FWD), which compare whole sets of fills, place averaged fills 16 and 186 times as far as random people but barely distinguish the right person from a random one.
No metric we tested does both.
Squared error is known to favor the average; the ladder measures by how much on real anatomy.
The mean of the 16 unrelated people whose visible surroundings most closely match the target's reaches 0.29 under MSE, averaged over the nine sets: closer than the opposite side of the same body in every set, and closer than a second photograph in faces.
This retrieval baseline weakens as the hidden region grows, so it must be reported at a stated size.
Pretrained inpainting models can then be read against these references: in faces, LaMa reaches \mfLamaHqImdbMse{} under MSE, closer than the person's own second photograph and only \mfLamaHqImdbAheadNearest{} ahead of that retrieval baseline.
A challenge-winning brain model reaches \mfBratsQtimMse{}, ahead of the baseline but behind the person's second scan.
The ladder is a deterministic calibration scale for inpainting metrics, and its references should be reported next to model scores.
\end{abstract}
